\subsection{From the Selmer ranks to the determinant locus}
\label{sec:patching-cotangent}
\hypertarget{source-ordinary-patching-cotangent}{}

The matrix realization of Section~\ref{sec:patching-matrices} lets us identify
infinitesimal deformations. We now use their dimensions to compare primes in the
completed localization at a comparison point. Write
\(B\) for that local ring, \(F=k((t))\) for its residue
field and \(r_F:B\to F\) for the residue map. The scalar
section is a complete coefficient DVR \(\mathcal O\)
mapping to \(B\) and inducing an isomorphism on residue
fields. The realization over the relative coefficient ring has the following properties. The coefficient algebra of
the matrices is finite over \(B\), with underlying module
\(B\oplus M\). The entries in the diagonal blocks are in
\(B\), those in the off-block rectangles in \(M\).
The representation matrices span all matrices with the specified pattern over \(B\).

In the product the off-block terms have coefficients in
\[
 M\cdot M\subseteq\mathfrak m_B\subset B .
\]
Write \(I_M\) for the ideal generated by these products.
In the one-block case \(M=0\).
The matrices satisfy the ordered local flag product
identity and the polarization and inertia relations.
The underlying global representation specializes to the
curve point. We use the finite matrix lattice over the
unlocalized determinant ring, continuous in its maximal
ideal topology, and the scalar lifts from the curve
ring \(k[[t]]\) in the coefficient section. These are
the matrix and coefficient inputs of the calculation.
\label{cor:ordinary-patching-cotangent-1}

\paragraph{The limit cohomology.}
The estimates available on the curve are finite-depth estimates. The first step is to
identify what they bound after passage to the limit.

\begin{lemma}[Passing the finite-depth bounds to the marked limit]\label{lem:marked-limit-cohomology}
Use the fixed curve and the marked system of
Proposition~\ref{prop:marked-profinite-system}.
Let \(Z_{\rm ord,b}\) be the bounded first-order cocycles satisfying the
ordinary local conditions on the limit and \(B_{\rm b}\) their subspace of
bounded boundaries. Then
\[
       \dim_F(Z_{\rm ord,b}/B_{\rm b})\le D_0=nd+nq .
\]
For \(g+1\) even tuples on the limit there is a nonzero
\(F\)-linear combination admitting the paired odd lift.
Any such lift has zero odd value on the marked tame generators,
and its restriction in the odd directions is read in the base
space with the chosen basis \((\kappa_i)_{i=1}^g\).
Here bounded means that a common power of \(t\) makes the
cocycles and the matrices imposing their local conditions integral.
\end{lemma}

\begin{proof}
Use the marked limit of the auxiliary groups from
Section~\ref{sec:patching-marking}. The
marking gives surjections from one finitely generated
profinite group to the groups at the auxiliary levels
and at the limit. In each finite quotient of the marking
group the level kernels have the same images as the
limit kernel on an ultrafilter extending the tails.
The marking also identifies the coefficient action,
complex conjugation, and the local restriction maps
on any fixed finite quotient on that ultrafilter.

\proofstep{Obtaining the finite-level cochains from the marked limit.}
At depth \(m\) pull back the cocycles to the positive
subgroup of the marking group. Let \(d\) be the number of topological generators and \(H\) the finite semidirect product at the chosen depth. Evaluation on those generators gives
\[
 \#\operatorname{Hom}_{\rm cont}(G,H)\le |H|^d,
\]
proving finiteness for the affine representations considered here. Each such
cocycle factors through a finite quotient of the
marking group. Take the simultaneous matching event
for these quotients and the coefficient action.
Descent to one level is equivalent there to descent
to the limit. The local cocycle values agree under
the marked local maps. The same matrix imposes the ordinary condition on either side, and unramifiedness is preserved as well.
The matrices generating boundaries also give exactly
the same pulled-back boundaries.

For this depth \(m\), let \(E_m\) be the matching
event for all the required finite quotients. If
\(N_{m,1}\) and \(N_{m,2}\) are thresholds for the two
finite-level estimates, choose a level in
\[
                 E_m\cap\{N\ge\max(N_{m,1},N_{m,2})\}.
\]
It is nonempty because both sets belong to the ultrafilter.
The finite quotient bound and the lift with factor \(t^C\)
transfer by the matched equations. This choice may depend on
\(m\), while the multiplier \(t^C\) is fixed. \label{cor:ordinary-patching-cotangent-2}

Choose \(e\) integral first-order cocycles with independent
classes over \(F\). The equation that an \(A\)-linear
combination is the boundary of a skew matrix is a system
of linear equations in the coefficients and the entries of
that matrix. Its rows span a finite \(A\)-module, so one
finite matrix supplies the same system at every \(t^m\).
Let \(B\) be this matrix, with domain the direct sum of
the coefficient module \(A^e\) and the module of integral skew matrices,
and let \(p\) project onto \(A^e\). Over \(F\) a relation among the chosen classes has zero coefficient vector.

In terms of the projection, 
\[
 \ker(B\otimes F)\subseteq\ker(p\otimes F).
\]
 Thus \(p\otimes F\) factors through the image of \(B\otimes F\); multiplying the factor map by \(t^{c_0}\)
for some fixed \(c_0\ge0\) makes it integral. Equivalently,
the nonzero Smith divisors of \(B\) give such a common
clearing exponent. Clearing denominators in the factor map gives an integral matrix \(L\) satisfying
\[
 LB=t^{c_0}p .
\]
For a vector satisfying the equations modulo \(t^m\),
\[
 Bv=t^mw\quad\Longrightarrow\quad
 t^{c_0}p(v)=L Bv=t^m Lw\equiv0\pmod{t^m}.
\]
Thus the kernel of the map from \((A/t^m)^e\) into the
finite-depth quotient is contained in
\((t^{m-c_0}A/t^m)^e\) for \(m\ge c_0\), giving an image
of size at least \(Q^{e(m-c_0)}\).
Proposition~\ref{prop:marked-first-order-bound} bounds that
quotient by \(Q^{mD_0+C}\) after finite-depth comparison.
Taking logarithms to base \(Q\) in these two estimates, then dividing by \(m\), gives
\[
 e(m-c_0)\le mD_0+C\quad\Longrightarrow\quad e\le D_0\qquad(m\to\infty).
\]

For each \(m\) apply the obstruction bound of
Proposition~\ref{prop:marked-odd-obstruction-bound} to the
\(g+1\) tuples reduced modulo \(t^m\). This is a bound for
lifting cocycles together with matrices imposing ordinarity. The failure
of a lifted tuple to satisfy the cocycle equation is an odd
two-cocycle, and its local corrections are the trivializations
in \textup{(11)}. Changing either lift adds the corresponding
coboundary in that complex. Therefore a unit-coefficient
combination multiplied by the fixed \(t^C\) has a lift.
Only the resulting one-cocycle and its local matrices are
transferred between the finite groups, through the finite
affine representations just described.

At each fixed depth there are finitely many coefficient vectors,
lifted cocycles and local matrices. A diagonal subsequence gives a
compatible reduction at every depth, hence an integral lift on
the limit. Write \((u_i)\) for the compatible coefficient vector. Unit-primitivity survives passage to the limit: 
\[
 \min_i v_t(u_i)=0,\qquad (t^Cu_i)_i\ne0\quad\text{in }F^{g+1}.
\]
 This gives the required nonzero combination. For a tame generator \(\tau_j\) the even data factor
through the base, where its image is trivial. The relation
\([\phi_j,\tau_j]=1\) therefore gives
\[
  \operatorname{Ad}(r(\phi_j))\,z(\tau_j)=z(\tau_j)
\]
for its odd cocycle value. The two Frobenius blocks have disjoint
spectra, so their off-diagonal centralizer is zero. The cocycle
rule and continuity now kill all of the closed normal subgroup
generated by the \(\tau_j\). Consequently the lift descends in
these directions to the ordinary odd space of the base, with
basis \((\kappa_i)\).
\label{cor:ordinary-patching-cotangent-3}

\proofstep{Passing to the fraction field with bounded denominators.}
Over \(F\) the three conclusions apply to bounded
continuous cocycles, meaning that some power of
\(t\) makes the cocycle integral and continuous on
every truncation. Clear one power for all entries
and for the finitely many local matrices, then
apply the integral conclusion. The cochain identities and local matrix equations are linear over the
coefficient ring and that ring injects into \(F\).

\label{cor:ordinary-patching-cotangent-4}

\end{proof}

Put
\[
       J_2=\mathfrak m_{\mathcal O}B+\mathfrak m_B^2+I_M,
       \qquad T=\bigl(\mathfrak m_B/J_2\bigr)^\vee .
\]
An element \(\alpha\in T\) extends uniquely to a derivation
\(B\to F\) at \(r_F\) killing \(\mathcal O\) and \(I_M\).
Indeed every element of \(B\) differs from a lift from
\(\mathcal O\) by an element of \(\mathfrak m_B\) and
products of two elements in \(\mathfrak m_B\) give zero.
The same equivalence identifies
\(\bigl(M/\mathfrak m_B M\bigr)^\vee\)
with additive maps \(\mu:M\to F\) that satisfy
\(\mu(bm)=r_F(b)\mu(m)\).
Apply \(\alpha\) to the diagonal block entries and
\(\mu\) to the off-block entries of the representation
matrices and multiply each result on the right by
the inverse of the base matrix. Denote the resulting
cocycles by \(z_\alpha\) and \(z_\mu\).
They are even and odd, respectively.

\begin{proposition}[Detection of the even and odd cotangent spaces]\label{prop:marked-cotangent-detection}
Use the normalized completed matrix model of
Proposition~\ref{prop:marked-finite-ch-algebra}. The even
derivation \(\alpha\in T\) and the odd functional
\(\mu\in(M/\mathfrak m_BM)^\vee\) give bounded cocycles satisfying
the indicated local and global conditions. Passing to cohomology detects each input separately: 
\[
 [z_\alpha]=0\ \Longrightarrow\ \alpha=0,\qquad [z_\mu]=0\ \Longrightarrow\ \mu=0.
\]
 In the two-block case choose
\(m_j\in M\) representing the coefficient of \([\kappa_j]\)
in the odd detection map. Then \(m_1,\ldots,m_g\) generate \(M\).
\end{proposition}

\begin{proof}
\proofstep{The paired test with its local relations.}
To see the multiplicative and local assertions for
these cocycles, use the paired test algebra over \(F\) with nilpotent basis
\(a_i,b_i,v\), products \(a_i b_j=\delta_{ij}v\) and all other
products of two basis elements zero. Indices range over whichever
finite family is being tested. For the original odd basis,
\(r_1=(1+\sum_j a_j\kappa_j)r\). Its adjoint on the square-zero
ideal with even \(b_j\) and odd \(v\) coordinates is the action
in the lifting assertion. Suppose that
\(\Phi:M\to F\) is \(\mathcal O\)-linear via \(r_F\)
and kills \(J_2M\). For \(m\in M\) define
\[
       d_m(b)=\Phi(bm)-r_F(b)\Phi(m).
\]
The map \(d_m\) kills \(\mathcal O\) and \(I_M\).
It is a derivation by writing
\(b=b_{\mathcal O}+b_1\) with \(b_1\in\mathfrak m_B\)
and using \(\Phi(\mathfrak m_B^2M)=0\).

Moreover \(d_{bm}=r_F(b)d_m\).
In a basis \(\alpha_i\) of \(T\) write
\(d_m=\sum_i\mu_i(m)\alpha_i\).
Each \(\mu_i\) is linear via \(r_F\).
Now the map on the coefficient algebra,
\[
 b+m\longmapsto
 r_F(b)+\sum_i a_i\mu_i(m)+\sum_i b_i\alpha_i(b)+v\Phi(m),
\]
is a ring map. In a mixed product \(bm\), the defining
derivation identity gives
\[
 \Phi(bm)=r_F(b)\Phi(m)+d_m(b)
         =r_F(b)\Phi(m)+\sum_i\alpha_i(b)\mu_i(m).
\]
Bilinearity and the multiplication table recover the required coefficient without mixed terms:
\[
 \left(\sum_i x_i a_i\right)\left(\sum_j y_j b_j\right)=\sum_{i,j}x_i y_j\delta_{ij}v=\left(\sum_i x_i y_i\right)v .
\] Products of
two odd elements have zero image: in the source they belong
to \(I_M\), killed by \(r_F\) and each \(\alpha_i\), and in
the target the odd-odd products are zero.
The same map works with any family \(\mu_i,\alpha_i\)
satisfying this formula for \(d_m\). By setting all
odd or all even data to zero, one obtains the respective
individual square-zero deformations defining
\(z_\alpha\) and \(z_\mu\).  \label{cor:ordinary-patching-cotangent-5}

The tests are bounded in the required topology.
Here is the role of the finite lattice in this argument.
Denote the nilpotent ideal of the test ring by \(J\) and choose \(r\) with \(J^r=0\). For any \(x_i\in J\) one has 
\[
 x_1^{e_1}\cdots x_d^{e_d}=0\qquad(e_1+\cdots+e_d\ge r).
\]
 The image of the unlocalized
determinant ring is contained in a finite
\(k[[t]]\)-algebra in that target. Indeed the scalar lifts
give the whole curve ring as constants. The kernel
of specialization to the curve ring in the unlocalized
determinant ring is finitely generated. The images
of those generators are nilpotent. If \(x_1,\ldots,x_d\) are the generator images, the scalar span of the following finite list already contains their image algebra:
\[
 \{x_1^{a_1}\cdots x_d^{a_d}:a_i\ge0,\ a_1+\cdots+a_d<r\}.
\]

A finite module of matrix
entries and differences follows by taking the images
of the generators of the original matrix lattice.
The images of the entries reduce to integral scalars.
Adjoining these entries using their nilpotent
differences from the scalars again gives a finite
algebra. If a determinant-ring element lies in the
\((j+r)\)-th power of its maximal ideal, expand its
products of maximal ideal elements into scalar parts
and nilpotent parts. The scalar parts lie in \((t)\)
and any term with at least \(r\) nilpotent factors
vanishes. The image is thus inside \(t^j\) times
the finite algebra. Continuity of the original
group matrices in their matrix lattice gives
continuity in these lattices for the test. Any
\(F\)-linear coordinate functional on them requires
only one fixed denominator.
\label{cor:ordinary-patching-cotangent-6}

\proofstep{The ordered local eigenspaces in the paired test.}
For a representation and ordered characters over the
test ring the ordinary relation is
\[
              \prod_{i=1}^n(\rho(h_i)-\psi_i(h_i))=0
                  \qquad(h_i\in G_v).
\]
The factors occur in order from \(1\) to \(n\).
By multilinearity the same identity holds on the group
algebra. Choose a finite linear combination \(z\) of group
elements for which the reduced character values at \(z\)
are pairwise different. Such a combination exists over
the infinite field \(F\): for each pair some group
element distinguishes it, and finitely many proper
linear hyperplanes cannot cover the space spanned by
these elements. The character differences lift to
units. Setting every factor at \(z\) gives the split
polynomial identity on \(\rho(z)\). The Lagrange
idempotents \(p_i\) of this polynomial are orthogonal
and sum to one.

For \(i\ge j\), put \(h\) in slot \(k\) with
\(j\le k\le i\), and \(z\) in all other slots.
Projecting the ordered product to \(p_i\) on the left
and \(p_j\) on the right gives
\[
 \left(\prod_{a<k}(\psi_i(z)-\psi_a(z))\right)
     p_i(\rho(h)-\psi_k(h))p_j
 \left(\prod_{a>k}(\psi_j(z)-\psi_a(z))\right)=0.
\]
The first product omits \(i\), the last omits \(j\);
both are units. If \(i>j\), the scalar term has
\(p_i p_j=0\), giving upper triangularity. If \(i=j=k\),
the surviving diagonal entry is \(\psi_i(h)p_i\).

\proofstep{Transport of the ordered relations through a nilpotent change.}
For two such systems agreeing modulo a square-zero
ideal, their idempotents \(p_i,p_i'\) agree modulo
that ideal. The matrix \(D=\sum_i p_i'p_i\) reduces to \(1\).
Orthogonality in each of the two idempotent systems gives
\[
       p'_jD=p'_jp_j=Dp_j,\qquad D^{-1}=1-(D-1).
\]
The inverse formula uses the square-zero ideal. Thus \(D\) carries
each old eigenspace to its new counterpart over the same ring. In the original ordinary frame the
original projectors are upper triangular, and a
basis made of their images of the respective standard
basis vectors differs from that frame by an upper
triangular unipotent change. Therefore
\(D^{-1}\rho'(h)D\) is upper triangular in the original
frame with the new prescribed diagonal. Both projectors
and \(D\) belong to the matrix subalgebra of the
allowed coefficient pattern. This follows from their
polynomial expressions and from inversion of scalar
units by a finite series in the nilpotent ideal.

Write \(D=1-Y\) and \(\rho'(h)=(1+\delta(h))r_1(h)\), with
\(Y,\delta(h)\) in the square-zero ideal. Multiplying the three
matrices gives
\[
  D^{-1}\rho'(h)D\,r_1(h)^{-1}
     =1+\delta(h)-\bigl(r_1(h)Y r_1(h)^{-1}-Y\bigr).
\]
Thus conjugation subtracts the coboundary of \(Y\), leaving a matrix in the required ordinary subspace.
On the prescribed local torsion subgroup the characters have fixed value,
so the resulting diagonal cocycle there vanishes.

A square-zero difference \(z\) therefore satisfies the differentiated multiplication and polarization laws 
\[
 z(gh)=z(g)+\operatorname{Ad}(r_1(g))z(h),\qquad z(\sigma h)=-z(h)^{\mathsf T},
\]
 in the polarized frame. The inertia identity gives the unramified
condition away from \(\ell\). This proves the local and
global claims for \(z_\alpha,z_\mu\) and for the
square-zero difference from \(r_1\) when the test
representation has the same \(a_i\) coordinates.
\label{cor:ordinary-patching-cotangent-7}

\proofstep{Detection modulo boundaries.}
If \(z_\alpha\) is a boundary, conjugation by a matrix
congruent to one kills the first-order representation
change. The determinant coefficients are then constant.
The flag characters are also constant. The ordered product
identity and the separating element of the group algebra
just used recover each character by its own projector.
The auxiliary Frobenius roots are unchanged because
their reduced differences are nonzero and the roots
are uniquely lifted for the same polynomial. On the corresponding spectral projectors the tame matrix is a conjugate of the constant identity matrix.

In particular for every allowed change of frame \(C\),
\[
 C\,\mathrm{id}\,C^{-1}-\mathrm{id}=0,
\]
so these tame variables vanish. The scalar parameter is fixed by the coefficient
section. The determinant coefficients, flag characters,
auxiliary roots and tame variables, and the scalar
parameter topologically generate the unlocalized ring.
Agreement on these elements gives agreement on the ring
by the finite lattice argument above. Localization and
completion preserve this equality in the test ring.
If \(J\) is the nilpotent test ideal and \(J^r=0\), each local test map \(f\) satisfies 
\[
 f(\widehat{\mathfrak m}^{\,r})\subseteq J^r=0.
\]
 Agreement on the generators thus extends to the completion and forces \(\alpha=0\).  \label{cor:ordinary-patching-cotangent-8}

Suppose \(z_\mu\) is the boundary of \(Y\).
Undoing the multiplication by the inverse base matrix, the
odd functional on each representation matrix is its scalar
specialization commuted with \(Y\). Both sides are \(B\)-linear
via \(r_F\), so the identity also holds on their \(B\)-span.
By the matrix-model hypothesis this span contains every
off-block matrix with a single entry \(m\in M\). Its scalar
specialization is zero, so the identity gives \(\mu(m)=0\).
This proves odd detection.

Write the odd detection map as
\[
 [z_\mu]=\sum_{j=1}^g \mu(m_j)[\kappa_j].
\]
Such \(m_j\) exist because the coefficients in the fixed basis
are linear functionals on \((M/\mathfrak m_BM)^\vee\), and
this finite-dimensional vector space has the natural double
dual. If the reductions of all \(m_j\) failed to span, there
would be a nonzero functional vanishing on them, contradicting
detection. Apply Nakayama's lemma to get the asserted generators.
\label{cor:ordinary-patching-cotangent-9}

\end{proof}

\paragraph{The spaces of liftable directions.}
To recover information about individual relations, we first separate the directions
that lift together. Let \(U\subseteq T^g\) consist of the tuples
\((\alpha_j)\) whose even cocycles \((z_{\alpha_j})\)
lift in the module with even \(b_j\) and odd \(v\)
coordinates, for the adjoint action of \(r_1\).
It is a linear subspace, since the possible lifts
and their local correction matrices satisfy linear equations.
The obstruction test on the limit, formulated in
Lemma~\ref{lem:marked-limit-cohomology}, shows that any
\(g+1\) vectors in \(T^g/U\) are linearly dependent, hence
\(\operatorname{codim}U\le g\).

\begin{proposition}[Dimension constraints on the liftable directions]\label{prop:marked-independent-lift-directions}
Let \(U\subset T^g\) consist of tuples whose even cocycles
lift to the module with even \(b_j\) and odd \(v\) coordinates
for the adjoint action of \(r_1\). Put
\(W=\{\alpha:e_j\alpha\in U\text{ for all }j\}\). Then
\[
 \operatorname{codim}U\le g,\qquad \dim_F W+g^2\le D_0 .
\]
For each \(\Phi\in(M/J_2M)^\vee\), the derivative tuple
\((d_{m_j})_{j=1}^g\) of the paired test belongs to \(U\).
\end{proposition}

\begin{proof}
\proofstep{Testing liftability by functionals on the obstruction module.}
For each functional \(\Phi\) on \(M/J_2M\), use
the paired test above. Its odd cocycle in coordinate
\(a_i\) has class
\(\sum_j\mu_i(m_j)[\kappa_j]\). Conjugation by one plus
a sum of the \(a_i\) times skew odd boundary generators
arranges equality of cocycles in these coordinates.
This leaves the even \(b_i\) coordinates unchanged.
Write \(C_{ij}=\mu_i(m_j)\).
Replace the \(a_i\)-coordinates by \(g\) coordinates
whose prescribed cocycles are \(\kappa_j\), and
replace the \(b_i\)-coordinates by their transpose
contractions with \(C\). Keep the scalar and \(v\)
coordinates. The resulting entries define a
representation even if \(C\) is rectangular or singular.
For multiplication, a pair of odd coordinate vectors
under this change satisfies
\[
           (C a)\cdot b=a\cdot(C^t b),
\]
so the product has exactly the old mixed coefficient.

The constant and first-order coordinates multiply by
the same formulas as before. The same calculation on
the ordered products, the polarization, and the inertia
identities preserves each of them. In relations with
zero right-hand side the new \(a_j\) parts vanish from
the corresponding relation for \(r_1\), while the other
parts vanish by the contraction.  Divide the new
representation by \(r_1\) on the right and subtract
one. The difference lies in the square-zero ideal
\((b_j,v)\), hence it is a cocycle for the adjoint action
of \(r_1\). The local flag argument above applies.
Its even coordinates are precisely the cocycles
of the \(g\) derivatives \(d_{m_j}\).
Consequently
\[
                  (d_{m_j})_{j=1}^g\in U
                    \qquad\text{for all }\Phi\in(M/J_2M)^\vee .
\]
Here this dual uses the \(F\)-vector space structure
from \(\mathcal O\), including the \(\mathcal O\)-linearity
required of \(\Phi\).  \label{cor:ordinary-patching-cotangent-10}

\proofstep{The subspace of independently liftable coordinates.}
Let \(e_j\alpha\) be the tuple supported in position \(j\). The condition on all positions is the intersection of the inverse images
\[
 W=\bigcap_j e_j^{-1}(U)=\{\alpha\in T:e_j\alpha\in U\ \text{for all }j\}.
\] For each \(\alpha\in W\), choose a lift
\(\Psi_j\) of \(e_jz_\alpha\). Put \(w_j=[v]\Psi_j\), the mixed coefficient of the lift. With \(a_i b_j=\delta_{ij}v\), the candidate cochain and its test images are 
\[
 Z=z_\alpha+\sum_{i=1}^g a_iw_i,\qquad b_jZ=b_jz_\alpha+vw_j=\Psi_j.
\]
 The scalar part and all products by
\(b_j\) determine such a first-order matrix. Thus the
cocycle, polarization, and inertia identities follow
from the corresponding identities for \(z_\alpha\)
and the \(\Psi_j\).

For the local condition, compare the even coordinates
of a matrix imposing ordinarity on \(\Psi_j\) with the tuple supported in
position \(j\) whose entry imposes the same condition on \(z_\alpha\). If \(D_{ij}\) is the difference in position \(i\) and \(\mathfrak b\) denotes the upper triangular space, local compatibility says 
\[
 \operatorname{Ad}(r(h))D_{ij}-D_{ij}\in\mathfrak b\qquad(h\in D_v).
\]

Each \(D_{ij}\) is therefore upper triangular over \(F\),
by the lower-entry induction used for the local
invariants in Section~\ref{subsec:odd-obstruction}.
In the ordinary first-order frame subtract their
lifts from the \(b_i\) coordinates of the matrix for \(\Psi_j\). These lifts are upper triangular.
Their boundaries stay upper triangular with diagonal
zero and do not change the required local condition.
Let \(Y\) impose the ordinary condition on \(z_\alpha\). The corrected matrix \(Y_j\) satisfies \([b_i]Y_j=\delta_{ij}Y\), and the same multiplication rule gives the combined matrix 
\[
 Y_* =Y+\sum_i a_i[v]Y_i,\qquad b_jY_*=Y_j.
\]
  \label{cor:ordinary-patching-cotangent-11}

For a basis of \(W\) these first-order classes, together
with the \(g^2\) classes \(a_j\kappa_i\), are independent.
Indeed a relation modulo boundaries in the scalar
coordinates first gives a boundary among the
\(z_\alpha\). The even detector forces the scalar
coefficients to vanish. The scalar part of the
boundary generator then commutes with the full
block span and is skew, hence is zero. In the remaining coordinates any relation has the form 
\[
 \sum_i c_{ij}[\kappa_i]=0\quad (1\le j\le g).
\]
 Independence of the odd classes gives \(c_{ij}=0\) for all \(i,j\).
We have constructed \(\dim W+g^2\) independent bounded
first-order classes on the same limit. Apply the dimension
bound of Lemma~\ref{lem:marked-limit-cohomology}.
\label{cor:ordinary-patching-cotangent-12}

\end{proof}

\paragraph{Independent coefficients of the odd relations.}
The liftable directions now have a role in the presentation: they distinguish the
coefficients of the odd relations one coordinate at a time.

\begin{proposition}[A presentation with independent odd-relation coefficients]\label{prop:marked-cotangent-presentation}
Use the paired cotangent detector of
Proposition~\ref{prop:marked-cotangent-detection} and the
per-coordinate liftable space \(W\) of
Proposition~\ref{prop:marked-independent-lift-directions}.
The completed marked ring \(B\oplus M\) is a quotient of
\[
 P=\mathcal O[[X_1,\ldots,X_{D_0},Y_1,\ldots,Y_g]] .
\]
There is a \(g\)-by-\(g\) matrix \(Z\) in the maximal ideal
of \(P\) for which the kernel contains the column \(ZY\).
The \(g^2+g\) entries of \(Z,Y\) are independent in the
relative cotangent space of \(P\) over its residue field.
\end{proposition}

\begin{proof}
Choose a spanning family of \(\operatorname{codim}U\) rows of
\(U^\perp\) and pad it to \(g\) rows. Denote its entries by
\(\phi_{ij}\in T^\vee\). For a tangent vector \(\alpha\), the entries of this matrix test all the supported tuples:
\[
 \phi_{ij}(\alpha)=0\ (\text{all }i,j)\quad\Longleftrightarrow\quad e_j\alpha\in U\ (\text{all }j).
\] Thus their common
annihilator is \(W\). A complement to their span has
dimension \(\dim W\le D_0-g^2\).
Form a spanning list of \(T^\vee\) by first listing all
\(g^2\) entries in designated slots and then a basis of this
complement, padding the end to length \(D_0\). When an entry
is expressed by projection to its designated slot, the
\(g^2\) coefficient vectors in \(F^{D_0}\) are independent.
The entries themselves may be dependent in \(T^\vee\).
Lift the whole spanning list to \(x_1,\ldots,x_{D_0}\in
\mathfrak m_B\) using \(T^\vee=\mathfrak m_B/J_2\).

\proofstep{Independent coefficient lifts in the ambient ring.}
Lift each coefficient to \(\mathcal O\) and write
\(z_{ij}\in\mathfrak m_B\) for the resulting linear
combination of the \(x_a\). For every such \(\Phi\), vanishing of the row on the tuple \((d_{m_j})\) gives
\[
 \Phi\!\left(\sum_j z_{ij}m_j\right)=0 .
\]
The dual separates \(M/J_2M\); thus the relation
holds modulo \(J_2M\).
Express the error as a combination of the
\(m_j\) with coefficients in \(J_2\), and
subtract it from the row. This gives an exact
relation \(\sum_j z'_{ij}m_j=0\).

The image of \(I_M\) in \(B\oplus M\) lies in the
square of its maximal ideal, because it comes from
products of two odd elements there. Consequently
the corrections have zero relative cotangent
class in \(B\oplus M\).
The complete local algebra \(B\oplus M\) is a
quotient of
\[
               P=\mathcal O[[X_1,\ldots,X_{D_0},Y_1,\ldots,Y_g]]
\]
by sending \(X_a,Y_j\) to \(x_a,m_j\); the
surjectivity follows from completeness and these
generators modulo the square and coefficient ideal.
Lift each corrected relation coefficient to
\(P\) with the same linear term modulo the
coefficient uniformizer. The entries of the
resulting matrix \(Z\), together with the \(Y_j\),
are independent in the relative cotangent space
of \(P\). The kernel contains the column \(ZY\).
 \label{cor:ordinary-patching-cotangent-13}

\end{proof}

Use Proposition~\ref{prop:marked-cotangent-presentation}
in the local comparison for
Proposition~\ref{prop:marked-chain-support}. We now relate
the determinant locus of the matrix \(Z\) to prime
containment in the completed patch ring.

\paragraph{Comparison of primes.}
We can now use dimensions to compare the primes through which support may pass. At a
two-block point the allowed target primes also depend on the odd module.

\begin{proposition}[Prime containment on the local comparison]
\label{prop:marked-stable-prime-inclusion}
For primes \(Z_0,Q_0\) of the one-block normalized comparison ring \(B\),
\[
      \dim(B/Z_0)\ge D_0+1\quad\Longrightarrow\quad Z_0\subseteq Q_0 .
\]
On the two-block comparison, a stable prime \(Z_0\) of that
quotient dimension lies in each of the admissible stable or
thick target primes \(Q_0\). Use the odd module \(M\) and the paired relation
matrix of Proposition~\ref{prop:marked-cotangent-presentation}.
\end{proposition}

\begin{proof}
We prove the local implication needed to propagate
the support. The one-block case has \(M=0\).
The same cotangent generators make \(B\) a quotient
of the regular local ring in \(D_0\) variables over
\(\mathcal O\). A prime of \(B\) with quotient dimension
at least \(1+D_0\) has zero preimage there. It is therefore
contained in every prime of \(B\).

\proofstep{The two-block ideal and its stable locus.}
In the two-block case write \(J=(ZY)\) in \(P\) and
\(I=J+(\det Z)\). The ideal \(I\) has a unique least
prime. To prove this, for a fixed index \(c\) adjoin
a column \(\lambda\) of variables with \(\lambda_c=1\)
and impose
\[
              Z\lambda=0,\qquad
              Y_i=\lambda_iY_c\quad(i\ne c).
\]
At any prime above \(\mathfrak m_P\), the entries of
\(Z\) and \(Y\) are still part of a regular parameter
system in the localized polynomial ring. Indeed one
can append lifts of regular parameters of the polynomial
fiber to the original parameters.

The chart equations
are cotangent independent there, with pivots in
\(Z_{ic}\) and \(Y_i\) for \(i\ne c\).
The quotient by those equations is regular local,
hence a domain. Let \(K_c\) be the contraction of
the chart ideal from the localization at
\(\mathfrak m_P P[\lambda]\). It contains \(I\):
the chart gives \(Y=Y_c\lambda\). Since \(Z\lambda=0\), the
adjugate identity gives
\[
 0=\operatorname{adj}(Z) Z\lambda=(\det Z)\lambda,\qquad
 0=(\det Z)\lambda_c=\det Z .
\]

Under any local map from \(P\) into a local domain
with a chart solution, \(K_c\) maps to zero.
Indeed extend the map by the given \(\lambda\) and
localize at the preimage of the target maximal ideal.
The chart ideal remains prime in this localization
and lies inside the extension of
\(\mathfrak m_P P[\lambda]\). A multiplier outside
\(\mathfrak m_P P[\lambda]\) in the localization criterion for
\(K_c\) thus lies outside the chart ideal and can be
cancelled. The resulting element belongs to the
chart ideal and maps to zero. In the generic
localized chart all \(\lambda_i\) are units.
Since \(\lambda_d\) is invertible here, normalization in chart \(d\) is given by
\[
 \lambda' = \lambda_d^{-1}\lambda,\qquad \lambda'_d=1.
\]
This takes a solution in chart \(c\) to one in chart \(d\), proving \(K_d\subseteq K_c\).
Reversing the indices proves equality.

\proofstep{Dimension after specializing the relation matrix.}
For a prime \(Q'\supseteq I\), choose a valuation ring
of the fraction field of \(P/Q'\) dominating that local
domain. If the image of \(Y\) is nonzero, choose a
coordinate of least valuation and divide by it.
The resulting \(\lambda\) lies in the valuation ring
and gives a chart solution. If \(Y=0\) in the residue field, the determinant condition on the image of \(Z\) gives a projective solution to the kernel equation:
\[
 \det Z=0\quad\Longrightarrow\quad \exists\,\lambda\ne0\text{ with }Z\lambda=0 .
\] Normalize it in the
same way. In either case the common contraction
\(K_c\) lies in \(Q'\). This proves the least-prime
assertion.

The dimension of \(P/J\) is at most \(1+D_0\).
In fact adjoining the \(g^2\) equations
\(Z_{ij}=\delta_{ij}Y_i\) gives the same radical as
setting all entries of \(Z,Y\) equal to zero,
since \((ZY)_i\) becomes \(Y_i^2\).
Cotangent independence makes all \(g^2+g\) entries of
\(Z,Y\) part of regular parameters, so killing them leaves
dimension \(\dim P-g^2-g\). Write \(J'\) for the ideal
obtained by adjoining the \(g^2\) equations to \(J\).
The radical comparison and the dimension inequality give
\[
       \dim(P/J)\le\dim(P/J')+g^2=\dim P-g .
\]
This is the asserted \(1+D_0\) bound. A prime over \(J\)
with quotient dimension at least \(1+D_0\) avoiding
some \(Y_j\) must therefore be the least prime over \(I\).
It contains \(I\) by the adjugate identity, and
strict containment above the least prime would
lose dimension.

\proofstep{The comparison at the stable and thick primes.}
Now let \(Z_0\subset B\) be stable with quotient
dimension at least \(1+D_0\). Stability in the
two-block case implies \(I_M\not\subseteq Z_0\).
Lift \(Z_0\) to the finite algebra \(B\oplus M\)
and take its preimage in \(P\). Its quotient
dimension is preserved by integrality. Some
\(Y_j\) avoids this prime, since containing all
the odd generators would force it to contain
their products. Therefore the preimage is the
least prime over \(I\).

For the target prime \(Q_0\) the two sufficient conditions in the comparison read
\[
 I_M\not\subseteq Q_0\quad\text{(stable)},\qquad \operatorname{Ann}_B M\subseteq Q_0\quad\text{(thick)} .
\] The determinant of the
exact relation matrix \((z'_{ij})\) annihilates
\(M\), by the adjugate identity. It lies in
\(Q_0\) in the case with nonsemisimple residue algebra. In the stable case its
product with \(I_M\) is zero, so primality gives
the same inclusion. A prime over \(Q_0\) in
\(B\oplus M\) thus pulls back to a prime over \(I\)
in \(P\). It contains the preimage for \(Z_0\).
Contracting gives \(Z_0\subseteq Q_0\).
For \(g=0\), the odd module is zero by its generation
statement, so the two-block stability hypothesis
cannot occur.
They give the prime inclusion once the patched
support provides the dimension lower bound at
the indicated curve points. \label{cor:ordinary-patching-cotangent-14}
\end{proof}
